\subsection{Adaptive tree approximation and causal completion}
Binev and DeVore~\cite{binevdevore} study near-optimal adaptive approximation trees and linear-in-tree-size construction. A greedy frontier is therefore not, by itself, a new approximation paradigm. Static singular quantization is likewise established: Graf and Luschgy~\cite{cantorquant} determine optimal quantization of the Cantor distribution, and Zhu, Zhou and Sheng~\cite{markovquant} study exact quantization-error orders for separated Markov-type measures. The profile in Lemma~\ref{lem:frontiercomparison} deliberately avoids claiming a new stationary fractal-dimension formula. Its proof is supplied because its constants must hold uniformly over all legal observation policies.

Theorem~\ref{thm:autonomous} concerns a different joint constraint. It first optimizes the actual full-history acquisition problem, then compiles that policy at a fixed worst-case observation budget into a finite autonomous controller. The proof must simultaneously preserve the acquired mass, the depth constraint, and a hard bound on all internal and terminal machine nodes. A static optimal quantizer does not provide that controller, while a near-best tree approximation does not by itself compare the separately optimized policy-dependent acquired laws. The finite-horizon pruning and cardinality-stability steps establish this comparison under explicit strong separation. They do not prove fast optimal policy planning: the dynamic program can be exponential, and a general lower bound on planning complexity is not claimed.

The correlated example is not a general unknown-kernel learning theorem. Its novelty claim is limited to a raw verification of observation-dependent acquisition within the uniform autonomous-completion theorem, with a computed strict feedback advantage. The continuous noisy examples use the different continuation mechanism and are not claimed to satisfy separated refinement. These distinctions delimit theorem-level content without promoting either realization into the general statement.
